\documentclass[12pt]{report}
\usepackage[margin=1in]{geometry}
\usepackage{setspace}
\usepackage{titlesec}
\usepackage{tocbibind}

\setstretch{1.5}
\setlength{\parindent}{0.5in}
\setlength{\parskip}{0pt}

\titleformat{\chapter}[display]
  {\normalfont\bfseries\centering}
  {\chaptername\ \thechapter}{0.5em}{}

\begin{document}

% Title page
\begin{titlepage}
  \centering
  {\Large Preliminary Design Tool for Flying Wing Unmanned Aerial Systems\par}
  \vspace{1.25in}
  {\large Malik D. Henry\par}
  \vspace{0.6in}
  {\large North Carolina Agricultural and Technical State University\par}
  \vspace{0.4in}
  {\large A thesis submitted to the graduate faculty\par}
  {\large in partial fulfillment of the requirements for the degree of\par}
  \vspace{0.4in}
  {\large Master of Science\par}
  {\large Program: Mechanical Engineering\par}
  \vspace{0.4in}
  {\large Major Professor: John P. Kizito\par}
  \vfill
  {\large Greensboro, North Carolina\par}
  {\large 2026\par}
\end{titlepage}

% Approval page
\clearpage
\thispagestyle{empty}
\begin{center}
  This is to certify that the master's thesis of\\[0.15in]
  Malik D. Henry\\[0.15in]
  has met the thesis requirements of\\[0.15in]
  North Carolina Agricultural and Technical State University\\[0.15in]
  Greensboro, North Carolina\\[0.35in]
  2026
\end{center}

\vspace{0.4in}
\noindent Approved by:\\[0.45in]
\noindent John P. Kizito \hfill Major Professor\\[0.45in]
\noindent \rule{3.5in}{0.4pt} \hfill Committee Member\\[0.45in]
\noindent \rule{3.5in}{0.4pt} \hfill Committee Member\\[0.45in]
\noindent \rule{3.5in}{0.4pt} \hfill Department Chair\\[0.45in]
\noindent Dr. Clay S. Gloster, Jr. \hfill Dean, The Graduate College

% Copyright page
\clearpage
\thispagestyle{empty}
\begin{center}
  Copyright by\\[0.15in]
  Malik D. Henry\\[0.15in]
  2026
\end{center}

% Front matter
\pagenumbering{roman}

\chapter*{Abstract}
The preliminary design of an aircraft is the stage where the major design decisions are
refined and iterated into a cohesive overall configuration, establishing the final general layout
before transitioning into detailed design. This phase focuses on aerodynamic, structural, and
performance analysis while balancing the constraints and requirements until a final configuration
is reached. The goal of this project is to consolidate these analyses for flying wing aircraft into a
single framework that enables rapid iteration from early geometry choices to actionable insights
moving into the detail design phase.

\chapter*{Biographical Sketch}
Malik D. Henry is a graduate student in the Mechanical Engineering program at North Carolina
Agricultural and Technical State University.

\chapter*{Dedication}
Optional.

\chapter*{Acknowledgments}
Optional.

\tableofcontents
\listoffigures
\listoftables

% Main matter
\clearpage
\pagenumbering{arabic}

\chapter{Introduction}
Add the problem statement, motivation, and project objectives.

\chapter{Literature Review}
Summarize prior work on flying-wing design methods, multidisciplinary analysis workflows,
and related tools.

\chapter{Methodology}
Describe geometry modeling, aerodynamic analysis, mission simulation, propulsion modeling,
structural checks, and export workflows used in this thesis.

\chapter{Results and Discussion}
Present design studies, validation results, and interpretation of outcomes.

\chapter{Conclusion}
Summarize findings, contributions, limitations, and next steps.

\chapter*{References}
\addcontentsline{toc}{chapter}{References}
\begin{itemize}
  \item Add references here.
\end{itemize}

\appendix
\chapter{Appendix}
Add supplementary material here.

\end{document}
